\documentclass[11pt]{article}
\usepackage[T1]{fontenc}
\usepackage{lmodern,microtype}
\usepackage[margin=1in]{geometry}
\usepackage{amsmath,amssymb,amsthm,mathtools,booktabs,xcolor}
\usepackage[colorlinks=true,linkcolor=blue!45!black,citecolor=blue!45!black,urlcolor=blue!45!black]{hyperref}
\usepackage[nameinlink,noabbrev]{cleveref}
\setlength{\emergencystretch}{2em}
\newtheorem{theorem}{Theorem}[section]
\newtheorem{proposition}[theorem]{Proposition}
\newtheorem{corollary}[theorem]{Corollary}
\theoremstyle{remark}\newtheorem{remark}[theorem]{Remark}
\newcommand{\E}{\mathbb E}\newcommand{\PP}{\mathbb P}
\newcommand{\R}{\mathbb R}\newcommand{\N}{\mathbb N}
\newcommand{\dd}{\,\mathrm d}\newcommand{\ind}{\mathbf1}
\newcommand{\Pois}{\operatorname{Pois}}
\DeclareMathOperator{\Cov}{Cov}
\title{Additive Cumulants and Computable Fixed-Sample Calibration\\
\large A research continuation of the discovery-inversion and transfer drafts}
\author{Christopher D. Long}
\date{November 17, 2026\\Research note, version 0.2}
\begin{document}\maketitle
\begin{abstract}
We give explicit sufficient integrability conditions for the additive cumulant
identity underlying the discovery-inversion draft. We then compute the transfer
draft's dominated-deletion calibrators by polynomial products, including the
lost-species correction needed for two-sided novelty prediction. Retaining the
full-count term gives an explicit logarithmic threshold cost. Its common data
score also permits candidate-dependent optimization of the Poisson effort,
without a multiple-testing penalty and while preserving convexity in the
relaxed population measure. These are proof-level extensions and numerical
research prototypes; we do not claim prior-art clearance, shortest intervals,
or conditional prediction coverage.
\end{abstract}

\section{Status and notation}
This note develops the additive-transform and computation interfaces of the
current discovery and transfer drafts \cite{LongDiscovery,LongTransfer}.
The capture--recapture companion \cite{LongCapture} supplies a counting-measure
version for independent capture histories and exact conditional kernels for
fixed-size and inverse sampling. Classical ingredients are cited at use;
detailed positioning appears in \cref{sec:prior}.

Let $p_i>0$, $\sum_i p_i=1$, with finitely or countably many disjoint species,
and set $\mu_p=\sum_i p_i\delta_{p_i}$. At Poisson effort $t$, the independent
counts are $X_i\sim\Pois(tp_i)$ and
$\Phi_r=\sum_i\ind\{X_i=r\}$, $r\ge1$.
For a finite direction $\lambda\in\R^d$, put $\lambda_0=0$ and
$\lambda_j=0$ for $j>d$. Write
\[
 b_r^t(x)=e^{-tx}(tx)^r/r!,\qquad
 K_\lambda^t(\nu)=\int \kappa_\lambda^t(x)\,\dd\nu(x),\quad
 \kappa_\lambda^t(x)=x^{-1}\log\!\left(1+
 \sum_{r=1}^d b_r^t(x)(e^{\lambda_r}-1)\right).
\]
The continuous extension is $\kappa_\lambda^t(0)=t(e^{\lambda_1}-1)$,
with $\lambda_1=0$ when $d=0$. The relaxed class is
$\mathcal P([0,1])$ unless additional restrictions are stated.
An atom at zero is a relaxation device, not a literal zero-probability species.
Total richness $\int x^{-1}\dd\nu$ still requires additional control near zero.

\section{The master identity with sufficient hypotheses}
Let $W_i$ be independent with law $Q_{p_i}$ and let
$\phi(x,W)\in\R^h$ be measurable. Any assertion below involving a kernel on
a relaxed class additionally assumes that its integral there is defined;
the theorem itself concerns the actual measure $\mu_p$.

\begin{theorem}[Additive cumulant identity]\label{thm:master}
Write $\phi_i=\phi(p_i,W_i)$. For finitely many species, assume the
moment generating functions in the following display are finite.
For countably many species, sufficient conditions at a specified
$\theta\in\R^h$ are
\[
 \sum_i\E\|\phi_i\|<\infty,\qquad
 \E\exp\!\left(\sum_i|\theta^\top\phi_i|\right)<\infty.
\]
Then $Z=\sum_i\phi_i$ exists absolutely almost surely and
\begin{equation}\label{eq:master}
 \log\E e^{\theta^\top Z}
 =\sum_i\log\E_{p_i}e^{\theta^\top\phi(p_i,W)}
 =\int x^{-1}\log\E_x e^{\theta^\top\phi(x,W)}\,\dd\mu_p(x),
\end{equation}
where the log series is absolutely convergent.
If, for some $\epsilon>0$,
\[
 \E\exp\!\left(\epsilon\sum_i\|\phi_i\|\right)<\infty,
\]
then the identity holds in a neighborhood of zero, and every joint cumulant
is obtained by differentiating the kernel and integrating:
\[
 \operatorname{cum}(Z_{j_1},\ldots,Z_{j_r})
 =\int x^{-1}\operatorname{cum}_x
       (\phi_{j_1}(x,W),\ldots,\phi_{j_r}(x,W))\,\dd\mu_p(x).
\]
These are sufficient, not necessary, conditions.
\end{theorem}
\begin{proof}
Absolute summability follows from Tonelli. For finite partial sums,
independence gives the product identity. Their exponentials are dominated by
the assumed integrable exponential, so dominated convergence identifies its
limit with the moment generating function of $Z$.
Put $V_i=\theta^\top\phi_i$. Jensen and monotonicity give
\[
 -\E|V_i|\le \log\E e^{V_i}\le \log\E e^{|V_i|}.
\]
The sums of the right-hand bounds are bounded by
$\log\E e^{\sum_i|V_i|}$, using independence and monotone convergence.
This proves absolute convergence of the log series. Its integral form is the
definition of $\mu_p$.

For differentiation, exponential domination gives summability of
$\E(e^{c\|\phi_i\|}-1)$ for every $0<c\le\epsilon$, since
$\sum_i(e^{c\|\phi_i\|}-1)\le e^{c\sum_i\|\phi_i\|}-1$.
On a sufficiently small complex polydisc about zero, the individual moment
generating functions therefore differ from one by a summable sequence,
uniformly. Discarding finitely many factors if necessary makes these
differences less than $1/2$, so their analytic logarithms converge normally.
The remaining finite factors have nonzero moment generating functions in a
possibly smaller neighborhood of zero. Termwise differentiation is now valid.
\end{proof}

For finitely many fingerprint coordinates, missing mass
$M_t=\sum_i p_i\ind\{X_i=0\}$, and future novelty
$Y=\sum_i\ind\{X_i=0,X_i'>0\}$ at independent effort $u$,
the sum of the per-species norms is bounded by a constant times
$1+K+L$, where $K=\sum_iX_i\sim\Pois(t)$ and
$L=\sum_iX_i'\sim\Pois(u)$. Thus the neighborhood condition holds
without a minimum probability.
The joint kernels are
\begin{align*}
 \kappa_{\lambda,\eta}^{t,M}(x)
 &=x^{-1}\log\!\left(1+\sum_r b_r^t(x)(e^{\lambda_r}-1)
                 +e^{-tx}(e^{\eta x}-1)\right),\\
 \kappa_{\lambda,\eta}^{t,u,Y}(x)
 &=x^{-1}\log\!\left(1+\sum_r b_r^t(x)(e^{\lambda_r}-1)
                 +e^{-tx}(1-e^{-ux})(e^\eta-1)\right).
\end{align*}
Their zero-rate limits are respectively
$t(e^{\lambda_1}-1)+\eta$ and
$t(e^{\lambda_1}-1)+u(e^\eta-1)$.
Means, covariances, affine prediction risks, and higher cumulants now follow
by differentiation. Untruncated long-horizon Good--Toulmin coefficients
do not satisfy the required exponential integrability in general.
The companion's separate moment and truncation arguments remain necessary.

For fixed directions $(\lambda_j,\eta_j)$ and nonnegative deterministic
weights with $\sum_j w_j\le1$, the region
\[
 \left\{(\nu,z):
 \sum_j w_j\exp\!\left(\lambda_j^\top\Phi+\eta_jz
                    -\int\kappa_j\,\dd\nu\right)\le1/\alpha\right\}
\]
is convex and covers the actual pair with probability at least $1-\alpha$.
This is the existing joint-calibration theorem with explicit sufficient
hypotheses supplied above. Projection onto a scalar $z$ gives an interval,
possibly empty or with unattained endpoints. Coverage is joint under the
specified experiment; it is not conditional on the observed sample.
Fixed-time validity alone does not imply validity at arbitrary stopping times.

\section{Computing dominated deletion}
Let $S_n$ consist of $n$ iid single-label draws, and let the positive
observed counts be $X_1,\ldots,X_D$.
For an indexed subset $I\subseteq[n]$, let $S_I$ denote its sub-sample.
Set $\pi_t(k)=e^{-t}t^k/k!$.
The transfer draft defines
\begin{equation}\label{eq:H}
 H_{\lambda,n}^t(S_n)=
 \sum_{k=0}^n\frac{\pi_t(k)}{\binom nk}
      \sum_{|I|=k} e^{\lambda^\top\Phi(S_I)}.
\end{equation}
The e-value at candidate $\nu$ is
$H_{\lambda,n}^t(S_n)e^{-K_\lambda^t(\nu)}$.
The retained Poisson weights must not be renormalized.

\begin{proposition}[Polynomial deletion algorithm]\label{prop:poly}
Define
\[
 P_\lambda(z)=\prod_{i=1}^D
      \left(\sum_{j=0}^{X_i}\binom{X_i}{j}e^{\lambda_j}z^j\right).
\]
Then $[z^k]P_\lambda=\sum_{|I|=k}e^{\lambda^\top\Phi(S_I)}$.
Sequential dense convolution computes all coefficients using
$O(n^2)$ arithmetic operations and $O(n)$ storage.
\end{proposition}
\begin{proof}
Selecting $j$ of the $X_i$ occurrences of species $i$ can be done in
$\binom{X_i}{j}$ ways, and contributes $e^{\lambda_j}$ to the score.
Multiplication sums independently over all species.
If $s_i=\sum_{h<i}X_h$, the convolution work is bounded by a constant times
$\sum_i(s_i+1)(X_i+1)=O(n^2)$, since $D\le n$ for $n>0$.
The empty sample has polynomial one.
\end{proof}
This is an arithmetic complexity statement, not a bit complexity or
rounding-error guarantee. The implementation uses log-domain convolution;
exact rational weights $e^{\lambda_j}$ are available in the verifier.

\subsection{Two-sided prediction and the lost-species coordinate}
For $m$ future draws let $Y$ be the number of labels unseen in $S_n$
but seen in the future sample. For the ordinary-draw case $B=1$,
the transfer draft's pathwise novelty sandwich gives the score
\begin{equation}\label{eq:Hjoint}
 H_{\lambda,\eta;n,m}^{t,u}(S_n)=
 \sum_{k=0}^n\sum_{\ell=0}^m
 \frac{\pi_t(k)\pi_u(\ell)}{\binom nk}
 \sum_{|I|=k}\exp\!\left\{
 \lambda^\top\Phi(S_I)-\eta_+(m-\ell)
                         -\eta_-\min(d_I,\ell)\right\},
\end{equation}
where $d_I$ counts species observed in $S_n$ but absent from $S_I$,
$\eta_+=\max(\eta,0)$, and $\eta_-=\max(-\eta,0)$.
Multiplying by
$\exp\{\eta y-K_{\lambda,\eta}^{t,u,Y}(\nu)\}$ gives a valid joint e-value.

\begin{proposition}[Bivariate deletion algorithm]\label{prop:bivariate}
Let
\[
 A_{k,r}=[z^kw^r]\prod_{i=1}^D
 \left(w+\sum_{j=1}^{X_i}\binom{X_i}{j}e^{\lambda_j}z^j\right).
\]
Then $A_{k,r}=\sum_{|I|=k,d_I=r}e^{\lambda^\top\Phi(S_I)}$, and
\[
 H_{\lambda,\eta;n,m}^{t,u}
 =\sum_{k,r}\frac{\pi_t(k)A_{k,r}}{\binom nk}
       \sum_{\ell=0}^m\pi_u(\ell)
       e^{-\eta_+(m-\ell)-\eta_-\min(r,\ell)}.
\]
Dense convolution uses $O(Dn^2)$ arithmetic operations and $O(nD)$
storage, with an additional $O(mD)$ operations for the inner weights.
When $\eta\ge0$, the correction is independent of $r$, so the
univariate algorithm suffices, with an additional $O(m)$ operations.
For $D=0$ the formula is interpreted directly.
\end{proposition}
\begin{proof}
A species is lost precisely when its selected count is zero, which contributes
the factor $w$. Otherwise it contributes its positive count and score.
Coefficient extraction and the indicated finite sums give the identity.
At each convolution stage there are at most $D+1$ lost-species coordinates;
the preceding univariate work bound gives the stated bound.
\end{proof}
These formulas do not establish the same complexity for arbitrary overlapping
features or dependent batches. Their inputs are ordinary single-label counts.

\section{The statistical cost of retaining full counts}
Write $E^P_\nu(S_n)=\sum_jw_j
 e^{\lambda_j^\top\Phi(S_n)-K_{\lambda_j}^t(\nu)}$
for the Poisson-calibrated formula evaluated at a fixed sample.
It need not itself be an e-value under the fixed-size law.
Let $E^D_\nu$ be its dominated-deletion counterpart.

\begin{proposition}[Deterministic threshold comparison]\label{prop:cost}
For a common effort $t$ and fixed nonnegative weights,
\[
 E^D_\nu(S_n)\ge\pi_t(n)E^P_\nu(S_n).
\]
Consequently,
\[
 \{E^D_\nu\le1/\alpha\}
 \subseteq\{E^P_\nu\le1/(\alpha\pi_t(n))\}.
\]
At $t=n\ge1$, the additional log threshold is
\[
 -\log\pi_n(n)=n-n\log n+\log(n!)
 =\tfrac12\log(2\pi n)+\frac1{12n}+O(n^{-3}).
\]
For joint prediction with common efforts $(t,u)$ the factor is
$\pi_t(n)\pi_u(m)$, and the log costs add.
\end{proposition}
\begin{proof}
Retain $k=n$ in \eqref{eq:H}. For prediction also retain $\ell=m$
in \eqref{eq:Hjoint}; then $d_I=0$ and both sandwich corrections vanish.
All omitted terms are nonnegative. Stirling's expansion gives the formula.
\end{proof}
Thus no confidence-level penalty does not mean no power cost.
The comparison is a deterministic region inclusion relative to the stated
formula with an enlarged threshold. It is not a universal comparison of
power, interval width, or statistical efficiency across experiments.

\section{Effort optimization for a common fixed-count score}
The full-count term admits a useful optimization that the full deletion
statistic does not automatically admit.

\begin{theorem}[Common-score normalization]\label{thm:common}
Let $V$ be a random vector under a fixed-count experiment with count
vector $\boldsymbol n$. Suppose that for every effort
$\boldsymbol t\in T$, a Poisson experiment has the same conditional law
of $V$ given its count vector $\boldsymbol n$, and
\[
 \E_{\mu,\boldsymbol t}e^{\theta^\top V}
       =e^{K_{\theta,\boldsymbol t}(\mu)},\qquad
 c_{\boldsymbol t}=\PP_{\boldsymbol t}\{
                 \boldsymbol K=\boldsymbol n\}>0
\]
does not depend on the population. Here the vector $V$ can include
realized targets as well as observed statistics.
For each candidate measure define
\[
 B_\theta(\nu)=\inf_{\boldsymbol t\in T}
       \{K_{\theta,\boldsymbol t}(\nu)-\log c_{\boldsymbol t}\}.
\]
Assume this is finite on the candidate class and measurable, and that
$K_{\theta,\boldsymbol t}$ is affine in $\nu$. Then
\[
 e_\nu(V)=\exp\{\theta^\top V-B_\theta(\nu)\}
\]
is an e-value at the true population under the fixed-count experiment.
Moreover $B_\theta$ is concave and $e_\nu$ is jointly convex in the
candidate measure and any scalar target coordinates substituted into $V$.
A fixed mixture over directions therefore yields a convex,
single-level confidence region.
On $\mathcal P([0,1])$, a finite nonempty effort grid with finite continuous
kernels suffices for all the stated regularity conditions.
\end{theorem}
\begin{proof}
Positivity and conditioning imply, for every $\boldsymbol t$,
\[
 c_{\boldsymbol t}\,
 \E_{\mu,\boldsymbol n}e^{\theta^\top V}
 \le e^{K_{\theta,\boldsymbol t}(\mu)}.
\]
The expectation on the left apart from $c_{\boldsymbol t}$ is the same
for all efforts. Taking the infimum of its deterministic upper bounds yields
$\E_{\mu,\boldsymbol n}e^{\theta^\top V}\le e^{B_\theta(\mu)}$.
Dividing proves the e-value assertion. An infimum of affine functionals is
concave. Hence $\theta^\top V-B_\theta(\nu)$ is convex in the candidate
variables; exponentiation and fixed nonnegative mixtures preserve convexity.
Markov's inequality proves coverage.
\end{proof}

For one window, $c_t=\pi_t(n)$ and $V=(\Phi,M_n)$ or simply $\Phi$.
For two independent windows, $c_{t,u}=\pi_t(n)\pi_u(m)$ and
$V=(\Phi,Y)$, with the joint novelty kernel above.
Conditioning independent Poisson processes on their total counts gives the
corresponding iid sample laws; this is exactly the hypothesis needed here.
Optimization is candidate-dependent but does not select a data direction:
the common score is fixed throughout.

On a finite grid the optimized e-value is the maximum of the
full-count e-values over that grid, yet remains valid because these e-values
are deterministic multiples of the same random variable.
Arbitrary maxima of e-values are not valid.
In particular, maximizing $H_{\lambda,n}^t(S_n)e^{-K_\lambda^t(\nu)}$
over $t$ changes the data score and is not justified by this proof.
The optimized full-count and a fixed-effort full-deletion e-value may be
combined using preassigned mixture weights. A maximum of those two objects
has no such automatic justification.

For $p=(0.5,0.3,0.2)$, $n=4$, and $\lambda=(2,-1)$, a declared grid
$t=0.05,0.10,\ldots,10$ selects $t=5.55$. Compared with $t=n$,
the log e-value increases by approximately $0.1368547442493213$ for every
sample at this candidate population. This is a finite-grid calculation,
not a claim that $5.55$ is the continuous optimum or that the resulting
intervals are optimal.

\section{Verification and next questions}
The accompanying verifier compares both polynomial products with direct
enumeration of indexed subsets using rational score weights, then checks
fixed-count e-value expectations by finite-population enumeration.
It reports 2150 checks, maximum relative numerical discrepancy below
$10^{-15}$ on the coefficient checks, and no expectation above one in the
tested calibration cases. Numerical corroboration is separate from the
proofs, and this does not certify floating-point behavior for large samples.
The numerical implementation uses logarithms to reduce overflow and
underflow; continuum moment optimization still requires its own certificates.

The immediate statistical question is whether a fixed mixture of optimized
full-count and deletion calibration materially improves projected interval
widths at moderate sample sizes. This should be tested on matched targets,
sampling designs, and parameter restrictions. Wu--Yang's support estimation
uses a minimum-mass class \cite{WuYang}. Painsky's intervals address individual
unobserved probabilities \cite{Painsky}; the recent maximum-unseen-probability
work uses a Bernoulli-product model \cite{Colombi}. These cannot be treated as
interchangeable missing-mass or richness benchmarks.

The remaining unification has a precise boundary: Poisson independence
supplies additive cumulants, positive or dominated kernels transport
calibration, and signed kernels transport expectation identities but can
amplify variation. Fixed-size calibration therefore survives, although its
log moment generating function need not be linear in $\mu$.
Overlapping-feature higher cumulants require joint incidence information.
The exact tank game remains a finite decision problem and a dual-certificate
example. The capture--recapture companion \cite{LongCapture} now develops the
closed-population tag--retag experiments, including exact fixed-size and
preassigned-recapture intervals and separate stopping-valid martingales.
For independent histories with fixed catchability vectors $q_i$, it uses
$\nu=\sum_i\delta_{q_i}$ and
\[
 \log\E e^{\theta^\top\sum_i\phi(q_i,W_i)}
 =\int\log\E_q e^{\theta^\top\phi(q,W)}\,\dd\nu(q).
\]
Population size is the unknown total mass of this counting measure.
Conditioning homogeneous Bernoulli captures on their sample totals recovers
the hypergeometric kernel; unequal catchabilities do not cancel in this way.
Marginalizing random catchabilities requires a different transform from
conditioning on a fixed population. A sharp bounded-heterogeneity
inverse-coupon procedure remains a separate target.

\section{Relation to prior work}\label{sec:prior}
Cumulant additivity for independent summands is classical; the abundance-measure
notation exposes its linear dependence on the population.
Poisson process foundations are described in \cite{LastPenrose}.
Moment recovery predates this program, notably Harris \cite{Harris} and
Daley--Smith \cite{DaleySmith}; generalized Chebyshev-system geometry is treated
by Karlin--Studden and Pinelis \cite{KarlinStudden,Pinelis}.
Positive transfer and deficiency belong to classical comparison of experiments
\cite{Blackwell,Torgersen}. These facts are foundations, not novelty claims.

The capture--recapture link does not begin with the present framework:
Chao's heterogeneous-population moment method and the unified estimating-
function approach of Chao and collaborators are discussed in detail in
\cite{LongCapture}. The common-score optimization proved here exploits
conditioning and positivity; neither that proof nor the polynomial algorithm
establishes first appearance or global statistical optimality.

\begin{thebibliography}{99}\raggedright
\bibitem{LongCapture}
C. D. Long. \emph{Capture--Recapture in a Sampling-Kernel Framework:
Exact recovery, estimator corrections, transfer, and calibration}.
Research draft, version 0.1, November 17, 2026.
Repository file \texttt{manuscripts/capture\_recapture.tex}.

\bibitem{LongDiscovery} C. D. Long.
\emph{Shape-Constrained Discovery-Rate Inversion}.
Research draft, version 0.3, November 17, 2026.
Repository file \texttt{manuscripts/discovery\_rate\_inversion.tex}.
\bibitem{LongTransfer} C. D. Long.
\emph{Certified Transfer and Generation in Partial-Population Inference}.
Research draft, version 0.2, November 17, 2026.
Repository file \texttt{manuscripts/transfer\_generation.tex}.
\bibitem{LastPenrose} G. Last and M. Penrose.
\emph{Lectures on the Poisson Process}. Cambridge University Press, 2017.
\href{https://doi.org/10.1017/9781316104477}{doi:10.1017/9781316104477}.
\bibitem{Harris} B. Harris.
Determining bounds on integrals with applications to cataloging problems.
\emph{Ann. Math. Statist.} \textbf{30}(2) (1959), 521--548.
\href{https://doi.org/10.1214/aoms/1177706266}{doi:10.1214/aoms/1177706266}.
\bibitem{DaleySmith} T. Daley and A. D. Smith.
Better lower bounds for missing species: improved non-parametric moment-based
estimation for large experiments.
\href{https://arxiv.org/abs/1605.03294v3}{arXiv:1605.03294v3}, 2019 revision.
\bibitem{KarlinStudden} S. Karlin and W. J. Studden.
\emph{Tchebycheff Systems: With Applications in Analysis and Statistics}.
Interscience, 1966.
\bibitem{Pinelis} I. Pinelis.
Tchebycheff systems and extremal problems for generalized moments: a brief survey.
\href{https://arxiv.org/abs/1107.3493}{arXiv:1107.3493}, 2011.
\bibitem{Blackwell} D. Blackwell.
Equivalent comparisons of experiments.
\emph{Ann. Math. Statist.} \textbf{24}(2) (1953), 265--272.
\href{https://doi.org/10.1214/aoms/1177729032}{doi:10.1214/aoms/1177729032}.
\bibitem{Torgersen} E. Torgersen.
\emph{Comparison of Statistical Experiments}. Cambridge University Press, 1991.
\bibitem{WuYang} Y. Wu and P. Yang.
Chebyshev polynomials, moment matching, and optimal estimation of the unseen.
\emph{Ann. Statist.} \textbf{47}(2) (2019), 857--883.
\href{https://doi.org/10.1214/17-AOS1665}{doi:10.1214/17-AOS1665}.
\bibitem{Painsky} A. Painsky.
Confidence intervals for parameters of unobserved events.
\emph{J. Amer. Statist. Assoc.} \textbf{120}(549) (2025), 226--236.
\href{https://doi.org/10.1080/01621459.2024.2314318}{doi:10.1080/01621459.2024.2314318}.
\bibitem{Colombi} A. Colombi, M. Beraha, A. Painsky, and S. Favaro.
Confidence intervals for maximum unseen probabilities, with application to
sequential sampling design.
\href{https://arxiv.org/abs/2601.20320}{arXiv:2601.20320}, 2026.
\end{thebibliography}
\end{document}
